\chapter{Conclusiones y Trabajo Futuro}
\label{cap:conclusiones}

\section{Conclusiones}
El objetivo de esta investigación era conseguir una infraestructura en la que a través de modelos de ML poder replicar el movimiento de una tela en un motor de videojuegos a tiempo real con baja latencia. Se buscaba identificar las limitaciones al traer técnicas que todavía están en el ámbito academico al desarrollo indie. \newline

Se crearon simulaciones de telas con geometrías simples y movimientos dinámicos provocados por un objeto colisionador en Unity. Durante la ejecución se extrajeron los datos de interés de la tela y se convirtieron a formato CSV. A partir de estos datasets se configuró un modelo que fuera capaz de aprender el comportamiento de la tela, adoptando poses de reposo y deformaciones realizadas por el colisionador sin colapsar. En el comienzo de la investigación se plantearon seis características (que conformarían 13 features) sobre los vértices, que podían ser relevantes para capturar la esencia de la tela en la simulación: la posición, el SDF (distancia a superficie colisionadora más cercana), la velocidad, el gradiente de SDF o normales respecto al colisionador, las coordenadas UVs y la máxima distancia de movimiento desde su origen en la malla de cada vértice. \newline

Se creó una base de entrenamiento con Pytorch, con la que cargar y limpiar los datos, preparar las estructuras necesarias (Datasets y modelos), entrenar y probar los mismos, extraer diversas métricas y finalmente exportarlos en formato ONNX. Fundamentalmente el modelo consistía en una entrada de features de vértices normalizados y aplanados en un vector, que producía una serie de deltas de desplazamiento de cada vértice. Sobre esta base se implementaron diferentes modelos y técnicas pertinentes para la mejora de la inferencia:
\begin{itemize}
    \item \textbf{Modelos:} Se implementaron dos tipos de modelo, . Conocer el contexto histórico de las simulaciones es esencial para dotar al modelo de cierta percepción temporal y aumentar la coherencia de la simulación.


    \item \textbf{Reducción de features:} De las características mencionadas, las tres más útiles para recrear el movimiento fueron la posición, SDF y velocidad. Los modelos finales cuentan únicamente con ellas como input.
    \item \textbf{Unrolling:} A esta RNN se le aplicó la ténica de ``unrolling'', que permite aplicar el algoritmo de retropropagación a través del tiempo (BPTT). De esta manera, el modelo infiere una secuencia de predicciones, corrigiendo su error de manera más acertada. 
    \item \textbf{Noise:} Modificar los datos de entrenamiento aplicandoles ruido, ha sido clave para conseguir un modelo más robusto, que pueda recuperarse de pequeños errores en la inferencia sin colapsar por el error acumulado.
    \item \textbf{Grandes datasets:} En la medida de lo posible, los entrenamientos con grandes cantidades de datos han demostrado naturalmente preparar al modelo para una mayor variedad de escenarios.
\end{itemize}

Los modelos finales se cargaron de nuevo en Unity, a través del paquete Sentis. El componente ClothML creado consiguió sustituir el componente de Cloth de Unity, replicando en la geometría un movimiento en base a los desplazamientos inferidos a través del modelo cargado. \newline

Con el pipeline descrito se plantearon experimentos sobre las distintas características que resultaban prometedoras para la optimización o mejora del resultado visual: la comparación del modelo recurrente (RNN) frente a un MLP, el efecto del uso del unrolling, la velocidad como parámetro de entrada y la reducción del input mediante PCA fueron analizados. A nivel visual se concluye que la combinación de la recurrencia y el unrolling aportan de estabilidad y conservación de la estructura a través del tiempo en la simulación evaluada. El siguiente modelo que resulta más interesante es el PCA, que además de aportar igualmente la conservación de la geometría, supone una mejora del rendimiento en tiempo de inferencia. No obstante, no alcanza el grado de realismo observado en el primero. \newline

A través de las diferentes características, se consigue acotar una dirección de estudio válida para el contexto investigado. Se hacen unas pruebas de usuario, que confirman que el resultado obtenido se acerca al comportamiento de una tela real, sin llegar al nivel de naturalidad propio de esta. Se disciernen las técnicas que favorecen la reproducción de los aspectos más relevantes para este tipo de simulaciones, marcando las siguientes limitaciones y pasos para el trabajo a futuro. 

\subsection{Limitaciones}

El estudio resultante queda con las siguientes limitaciones:
\begin{itemize}
    \item \textbf{Dataset pequeño} en comparación con el estado del arte. Para el nivel de simplicidad de la tela empleado en la simulación, el dataset generado ha sido suficiente para generar un comportamiento fiable a nivel visual. Sin embargo se podrían producir resultados mejores, así como menos dependencia con el objeto colisionador, introduciendo más diversidad de movimiento y número de frames totales. Por limitaciones de tiempo, presupuesto y hardware, en esta investigación se ha usado una media de 5000 frames por entrenamiento, mientras que en otras investigaciones que emplean bancos de simulaciones para entrenar sus modelos llegan a usar hasta $10^{6}$ frames \cite{holden_subspace_2019}.
    \item \textbf{Topología y generalización}. El modelo puede generar buenos resultados cuando la tela empleada en el entreamiento y simulación final coinciden. No es capaz de generalizar a otras topologías o mallas dinámicas.
    \item \textbf{Complejidad}. Como se menciona en el apartado del Dataset, debido a la dificultad para generarlas en Unity, el modelo no ha sido probado con telas más complejas o pegadas al cuerpo como camisetas o globos con forma. Podría no resultar óptimo e incluso no producir resultados aceptables en la inferencia. 
    \item \textbf{Optimización}. Con objeto de producir los mejores resultados posibles en el modelo y en el aspecto visual, se ha enfocado la fuerza de trabajo en el entrenamiento y no se ha llevado a cabo una optimización exhaustiva de la aplicación del modelo en Unity. Detalles como podría ser reducir el cálculo de distancias a los colisionadores no han entrado dentro del plazo de desarollo disponible.
\end{itemize}

- La simulacion elegida ha resultado ser... controversial. Ya al final del proceso, se llegó a la obvs de que este tipo de simulaciones con tanto movimiento e inercia se dejan siempre a mano de los solvers, usandose la IA más bien para reducir el coste de simular la ropa de personajes secundarios, arrugas en herramientas para artistas, vueltas al reposo... Se presenta la inercia academicamente como un caso de estudio, pero aun no se usa de forma profesional para resolver este problema.
- en la academia no es obj general d los estudios
- las técnicas que sí existen (loss físicos, optim topologicas, ya mencionados) too much
- en cuanto al alcance técnico no se cuenta con presupuesto ni máquinas pros, sino con portátiles de potencia reducida
- por ende entrenamientos reducidos = dataset pequeño, aplicaciones sin licencia por lo q todo en Unity




\section{Trabajo Futuro}

El trabajo a futuro de este estudio se puede centrar en varios esfuerzos:
\begin{itemize}
    %\item Realizar \textbf{pruebas de usuario} para comprobar si los resultados obtenidos se perciben como naturales. La obtención de feedback por parte de usuarios de videojuegos podría confirmar qué aspectos de la tela se han conseguido replicar y dónde se deberían producir mejoras.
    \item En cuanto a la infraestructura estaría bn plantear una automatización completa del flujo de trabajo. Como se ha estado iterando sobre los entrenamientos no. Pero estaría bn, con un modelo definitivo generar la herramienta completa de recogida de datos, entrenamiento y aplicación de los resultados a la propia geometría.
    \item Trabajar en la escalabilidad de nuestro modelo y herramienta. Más cosas de reducción como el PCA. No ha estado en el punto de mira, al tener recursos limitados, pero es crucial para el funcionamiento y necesario trabajo a futuro y remodelacion.
    \item Losses físicos. No hemos centrado en esto pero Interesante tanto por la restriccion física del espacio y conservacion de la geometria,  o generalizar topologia, como para el tamaño de los datasets.
    \item \textbf{Extender los datasets} con más situaciones y probar otros \textbf{modelos más complejos}. Han quedado por probar modelos de interés como las CNNs, GNNS y redes encoder-decoder.
    \item Probar con \textbf{otros tipos de tela}, otros materiales, prendas pegadas al cuerpo, y potencialmente aplicar \textbf{técnicas híbridas} de animación, simulación e inteligencia artificial para conseguir mejores resultados. Un ejemplo sería combinar técnicas de skinning con ML, atando la tela a huesos y dejando la inercia y detalles como arrugas a la inteligencia artificial.
    \item Añadir la posibilidad de \textbf{extraer datasets de software especializado} como Maya o Blender, e incluso de bancos externos de interés. La obtención de mejores referencias habilitaría la prueba de geometrías más complejas y de mayor calidad.
\end{itemize}